\documentclass[12pt]{article}
\usepackage{resource/shortex}
\usepackage{mathtools,amssymb}
\usepackage{enumitem}
\usepackage{datetime}
\usepackage{fancyhdr}
\usepackage{ifthen}
\usepackage{pifont}
\usepackage{mathrsfs}
\usepackage[normalem]{ulem}
%\RequirePackage[usenames,dvipsnames]{color}


%% showing solutions 
\newboolean{showsols}
\setboolean{showsols}{false}
\newcommand{\showsolutions}{\setboolean{showsols}{true}}
\newlength{\solspace}
\setlength{\solspace}{16em}
\newcommand{\nosolspace}{\setlength{\solspace}{0em}}

\newcommand{\solution}[2]{\ifthenelse{\boolean{showsols}}{{\par\textcolor{red}{Rubric criterion: #1}\newline\textcolor{blue}{#2}}}{\vspace{\solspace}}}
\newcommand{\npforprint}{\ifthenelse{\boolean{showsols}}{}{\newpage}}
\newcommand{\vsforprint}[1]{\ifthenelse{\boolean{showsols}}{}{\vspace{#1}}}

%% grading
\newcommand{\grades}[1]{~\newline\noindent\framebox[\textwidth][c]{E\quad \quad \quad \quad S \quad \quad \quad \quad R \quad \quad \quad \quad N}}

\newcommand{\questiongrades}{\noindent\emph{Each item is marked $\checkmark$, $\checkmark\!-$, or \ding{55}:}\begin{center}\emph{$\checkmark$ = capable  $~~$ $\checkmark\!-$ = mostly capable  $~~$ \ding{55} = not yet}\end{center} }

\newcommand{\triAssessments}{$\checkmark\quad\checkmark\!-\quad$\ding{55}\quad}
\newcommand{\biAssessments}{$\checkmark\quad$\ding{55}\quad}

%% line for writing name and BUID
\newcommand{\nameBUID}{\noindent\textbf{Name:}\underline{\phantom{XXXXXXXXXXXXXXXXXXX}}\hfill\textbf{BUID:}\underline{\phantom{XXXXXXXXXXXXX}}\newline}


%% formatting 
%\allowdisplaybreaks[3]
\renewcommand{\headrulewidth}{.4mm} % header line width
\renewcommand{\arraystretch}{1.25}


%% page style 
\pagestyle{fancy}
\fancyhf{}
\rhead{Spring 2026}
\lhead{CAS MA 213: Basic Statistics and Probability}
%\rfoot{\thepage}

%--------------------
\begin{document}

\begin{center}
\textbf{\Large Project 1, Part 1 \\ Lab Activity}
\end{center}

\begin{center}
\begin{tabular}{ll}
\toprule
\textbf{Item} & \textbf{Date} \\
\midrule
Lab P1\_1 (today) & Fri, Sep 25 \\
\textbf{Deliverable 1 (outline) due} & \textbf{Thu, Oct 1, 9:00 PM} \\
Tutorial 4 due; Lab 4 in class & Fri, Oct 2 \\
\textbf{Deliverable 2 (slides + R code) due} & \textbf{Thu, Oct 8, 9:00 PM} \\
Tutorial 5 due; Lab 5 in class & Fri, Oct 9 \\
Lab P1\_2 in class & Fri, Oct 16 \\
\textbf{Deliverable 3 (recorded presentation) due} & \textbf{Thu, Oct 22, 9:00 PM} \\
\bottomrule
\end{tabular}
\end{center}

\subsection*{Goal for Today}
Get your group set up to work efficiently on Project 1 for the rest of the week. You are \textbf{not} expected to finish the Deliverable 1 outline in lab today --- the goal is a solid start: a data set, a draft research question, and a first pass at your analysis plans. You'll have until the Deliverable 1 deadline above to finish it together outside of lab. By the end of today's lab, every member should also know their role and the time of your next check-in.

\subsection*{Before You Start}
Copy the \href{https://docs.google.com/document/d/1u9wn9czn2f1G6CzEF46wkjXuLb_-Z1Dh_4oGpTxe7AE/edit}{P1\_D1\_Outline} Google doc and share it amongst your group members --- that is the document you will fill in as you go. This activity sheet just tells you what to do and when; it has no blanks of its own.

You have the rest of the lab time to work on this. If don't get to step 6 (assigning roles and scheduling your next meeting) with 10 minutes left, stop where you are and move straight to step 6 --- every group should leave lab today with roles assigned and a next meeting on the calendar.

\subsection*{Suggested Order}

\begin{enumerate}
    \item \textbf{Introduce group members ($\sim$3 min)}
    \begin{itemize}
        \item Go around the group, say your name and BU ID, and fill in outline Section 2 (Group Members).
        \item Identify your TF/LA mentor, and get their email address.
        \item \emph{Tip:} pick one person to be the ``scribe'' typing into the shared doc for this session --- it keeps the doc moving and avoids duplicate edits from four people at once.
    \end{itemize}

    \item \textbf{Explore data sets and discuss possible research questions ($\sim$10 min)}
    \begin{itemize}
        \item Browse candidate data sets (see the list of suggested sources below) and talk through what each one could support. 
        \item \emph{Tip:} favor a data set you can describe in one sentence (``one row = one flight'') over one that sounds impressive but is hard to explain. A clear observational unit makes every later step easier.
    \end{itemize}

    \item \textbf{Pick a data set and confirm it loads in R ($\sim$10 min)}
    \begin{itemize}
        \item Converge on one data set. Actually open the data in R/RStudio; don't just look at a preview online.
        \item Once it loads, fill in outline Section 5: the data source link/citation, what one row represents, how many rows there are, and what the important variables are (numerical vs.\ categorical).
        \item \emph{Tip:} if the import fails or the file is huge/messy, that is useful information \emph{now} --- better to discover it in lab than the night before Deliverable 1 is due. If it's not working, drop that data set and try your next candidate instead.
        \item \emph{Tip:} if the group is stuck between two options, pick the one with fewer missing-data or access problems --- you can always pivot early in the week, but not after Deliverable 2 is drafted.
    \end{itemize}

    \item \textbf{Draft your research question and analysis plans ($\sim$10 min)}
    \begin{itemize}
        \item Write your primary research question (outline Section 4), then draft the numerical and categorical analysis plans (Sections 6a/6b). For each, note what you expect to see.
        \item \emph{Tip:} state the question first, then ask ``is this comparative, descriptive, or about a relationship?'' If it sounds causal (``does X cause Y''), rewrite it as an association or comparison --- you can't establish causation from observational data.
    \end{itemize}

    \item \textbf{Go back and fill in the title, motivation, and data-quality check ($\sim$5 min)}
    \begin{itemize}
        \item Now that you know your data and question, it's much easier to write a good project title and motivation --- fill in Sections 1 and 3. Also note one data-quality issue to investigate (Section 6c).
        \item \emph{Tip:} your motivation doesn't need to be profound --- ``we're curious whether X'' or ``this affects people we know'' is a fine reason. Graders want to see genuine engagement, not a contrived justification.
    \end{itemize}

    \item \textbf{Assign roles by deliverable and plan your next meeting ($\sim$7 min)}
    \begin{itemize}
        \item Fill in the Roles by Deliverable grid (outline Section 7) and agree on a specific date/time (or async check-in method) for your next meeting, before Deliverable 1 is due.
    \end{itemize}
\end{enumerate}

\subsection*{Tip: Suggested Task List by Deliverable}

Use this as a menu when filling in the Roles by Deliverable grid.

\begin{description}[nosep, leftmargin=0pt, style=sameline, font=\normalfont]
    \item[\textbf{D1 --- Outline (write sections):}] 1--2 title/group, 3 motivation, 4 research question, 5 data source, 6a numerical plan, 6b categorical plan, 6c data quality
    \item[\textbf{D2 --- Plots/tables (R analysis):}] histogram, boxplot, scatter plot, numerical summary table, bar plot, categorical summary table
    \item[\textbf{D2 --- Slides:}] title, data description, numerical analysis, categorical analysis, conclusion \& future work, references
    \item[\textbf{D2 --- Other:}] R code cleanup and assembly, roles \& responsibilities summary
    \item[\textbf{D3 --- Sections:}] introduction, numerical analysis, categorical analysis, conclusion
    \item[\textbf{D3 --- Other:}] recording and video editing
\end{description}

\subsection*{Possible Data Sources to Consider (Not Limited to These)}
You can click the links below to explore possible data sources.
\begin{itemize}
    \item \href{https://www.openintro.org/data/}{OpenIntro Datasets}
    \item \href{https://cran.r-project.org/web/packages/fivethirtyeight/vignettes/fivethirtyeight.html}{fivethirtyeight R Package}
    \item \href{https://stat.ethz.ch/R-manual/R-devel/library/datasets/html/00Index.html}{datasets R package}
    \item \href{https://www.kaggle.com/datasets}{Kaggle datasets}
    \item \href{https://data.gov/}{US government's open data}
\end{itemize}

\end{document}
